% prezentace.tex -- Lekce 03: Podmínky
% preamble-prez.tex -- sdílená preamble pro prezentace (beamer)
% Includuje se z ucitel/lekceXX/prezentace.tex jako:
%   % preamble-prez.tex -- sdílená preamble pro prezentace (beamer)
% Includuje se z ucitel/lekceXX/prezentace.tex jako:
%   % preamble-prez.tex -- sdílená preamble pro prezentace (beamer)
% Includuje se z ucitel/lekceXX/prezentace.tex jako:
%   \input{../spolecne/preamble-prez.tex}

\documentclass[aspectratio=169]{beamer}

% --- Jazyk a font ---
\usepackage{polyglossia}
\setmainlanguage{czech}
\usepackage{fontspec}

% --- Téma ---
\usetheme{default}
\usecolortheme{default}

\definecolor{microbit-blue}{HTML}{1E4D8C}
\definecolor{microbit-lightblue}{HTML}{D6E4F7}
\definecolor{code-bg}{HTML}{F0F0F0}

\setbeamercolor{frametitle}{bg=microbit-blue, fg=white}
\setbeamercolor{title}{fg=microbit-blue}
\setbeamercolor{block title}{bg=microbit-blue, fg=white}
\setbeamercolor{block body}{bg=microbit-lightblue}
\setbeamercolor{itemize item}{fg=microbit-blue}
\setbeamercolor{itemize subitem}{fg=microbit-blue!70}
\setbeamerfont{frametitle}{size=\large, series=\bfseries}
\setbeamertemplate{navigation symbols}{}
\setbeamertemplate{footline}{%
  \leavevmode%
  \hbox{%
    \begin{beamercolorbox}[wd=.5\paperwidth,ht=2.5ex,dp=1.125ex,leftskip=6pt]{author in head/foot}%
      \usebeamerfont{author in head/foot}\textcolor{gray}{\small Úvod do Pythonu na Micro:bitu}
    \end{beamercolorbox}%
    \begin{beamercolorbox}[wd=.5\paperwidth,ht=2.5ex,dp=1.125ex,rightskip=6pt,right]{date in head/foot}%
      \usebeamerfont{date in head/foot}\textcolor{gray}{\small\insertframenumber\,/\,\inserttotalframenumber}
    \end{beamercolorbox}}%
}

% --- Česky korektní uvozovky ---
\usepackage{csquotes}
\newcommand{\uv}[1]{\enquote{#1}}

% --- Zdrojový kód (Python) ---
\usepackage{listings}
\lstset{
  language=Python,
  basicstyle=\ttfamily\small,
  backgroundcolor=\color{code-bg},
  frame=single,
  framerule=0pt,
  xleftmargin=6pt,
  xrightmargin=6pt,
  breaklines=true,
  showstringspaces=false,
  keywordstyle=\color{microbit-blue}\bfseries,
  stringstyle=\color{teal},
  commentstyle=\color{gray}\itshape,
  literate=
    {á}{{\'a}}1 {é}{{\'e}}1 {í}{{\'i}}1 {ó}{{\'o}}1 {ú}{{\'u}}1
    {ů}{{\r{u}}}1 {č}{{\v{c}}}1 {š}{{\v{s}}}1 {ž}{{\v{z}}}1
    {Á}{{\'A}}1 {É}{{\'E}}1 {Í}{{\'I}}1 {Ó}{{\'O}}1 {Ú}{{\'U}}1
    {Č}{{\v{C}}}1 {Š}{{\v{S}}}1 {Ž}{{\v{Z}}}1 {ň}{{\v{n}}}1 {ř}{{\v{r}}}1,
}

% --- Zvýraznění pojmů ---
\newcommand{\pojem}[1]{\textbf{\textcolor{microbit-blue}{#1}}}
\newcommand{\pyinline}[1]{\texttt{#1}}

% --- Grafika a diagramy ---
\usepackage{tikz}

% --- Ostatní ---
\usepackage{graphicx}
\usepackage{hyperref}
\hypersetup{colorlinks=true, urlcolor=microbit-blue}


\documentclass[aspectratio=169]{beamer}

% --- Jazyk a font ---
\usepackage{polyglossia}
\setmainlanguage{czech}
\usepackage{fontspec}

% --- Téma ---
\usetheme{default}
\usecolortheme{default}

\definecolor{microbit-blue}{HTML}{1E4D8C}
\definecolor{microbit-lightblue}{HTML}{D6E4F7}
\definecolor{code-bg}{HTML}{F0F0F0}

\setbeamercolor{frametitle}{bg=microbit-blue, fg=white}
\setbeamercolor{title}{fg=microbit-blue}
\setbeamercolor{block title}{bg=microbit-blue, fg=white}
\setbeamercolor{block body}{bg=microbit-lightblue}
\setbeamercolor{itemize item}{fg=microbit-blue}
\setbeamercolor{itemize subitem}{fg=microbit-blue!70}
\setbeamerfont{frametitle}{size=\large, series=\bfseries}
\setbeamertemplate{navigation symbols}{}
\setbeamertemplate{footline}{%
  \leavevmode%
  \hbox{%
    \begin{beamercolorbox}[wd=.5\paperwidth,ht=2.5ex,dp=1.125ex,leftskip=6pt]{author in head/foot}%
      \usebeamerfont{author in head/foot}\textcolor{gray}{\small Úvod do Pythonu na Micro:bitu}
    \end{beamercolorbox}%
    \begin{beamercolorbox}[wd=.5\paperwidth,ht=2.5ex,dp=1.125ex,rightskip=6pt,right]{date in head/foot}%
      \usebeamerfont{date in head/foot}\textcolor{gray}{\small\insertframenumber\,/\,\inserttotalframenumber}
    \end{beamercolorbox}}%
}

% --- Česky korektní uvozovky ---
\usepackage{csquotes}
\newcommand{\uv}[1]{\enquote{#1}}

% --- Zdrojový kód (Python) ---
\usepackage{listings}
\lstset{
  language=Python,
  basicstyle=\ttfamily\small,
  backgroundcolor=\color{code-bg},
  frame=single,
  framerule=0pt,
  xleftmargin=6pt,
  xrightmargin=6pt,
  breaklines=true,
  showstringspaces=false,
  keywordstyle=\color{microbit-blue}\bfseries,
  stringstyle=\color{teal},
  commentstyle=\color{gray}\itshape,
  literate=
    {á}{{\'a}}1 {é}{{\'e}}1 {í}{{\'i}}1 {ó}{{\'o}}1 {ú}{{\'u}}1
    {ů}{{\r{u}}}1 {č}{{\v{c}}}1 {š}{{\v{s}}}1 {ž}{{\v{z}}}1
    {Á}{{\'A}}1 {É}{{\'E}}1 {Í}{{\'I}}1 {Ó}{{\'O}}1 {Ú}{{\'U}}1
    {Č}{{\v{C}}}1 {Š}{{\v{S}}}1 {Ž}{{\v{Z}}}1 {ň}{{\v{n}}}1 {ř}{{\v{r}}}1,
}

% --- Zvýraznění pojmů ---
\newcommand{\pojem}[1]{\textbf{\textcolor{microbit-blue}{#1}}}
\newcommand{\pyinline}[1]{\texttt{#1}}

% --- Grafika a diagramy ---
\usepackage{tikz}

% --- Ostatní ---
\usepackage{graphicx}
\usepackage{hyperref}
\hypersetup{colorlinks=true, urlcolor=microbit-blue}


\documentclass[aspectratio=169]{beamer}

% --- Jazyk a font ---
\usepackage{polyglossia}
\setmainlanguage{czech}
\usepackage{fontspec}

% --- Téma ---
\usetheme{default}
\usecolortheme{default}

\definecolor{microbit-blue}{HTML}{1E4D8C}
\definecolor{microbit-lightblue}{HTML}{D6E4F7}
\definecolor{code-bg}{HTML}{F0F0F0}

\setbeamercolor{frametitle}{bg=microbit-blue, fg=white}
\setbeamercolor{title}{fg=microbit-blue}
\setbeamercolor{block title}{bg=microbit-blue, fg=white}
\setbeamercolor{block body}{bg=microbit-lightblue}
\setbeamercolor{itemize item}{fg=microbit-blue}
\setbeamercolor{itemize subitem}{fg=microbit-blue!70}
\setbeamerfont{frametitle}{size=\large, series=\bfseries}
\setbeamertemplate{navigation symbols}{}
\setbeamertemplate{footline}{%
  \leavevmode%
  \hbox{%
    \begin{beamercolorbox}[wd=.5\paperwidth,ht=2.5ex,dp=1.125ex,leftskip=6pt]{author in head/foot}%
      \usebeamerfont{author in head/foot}\textcolor{gray}{\small Úvod do Pythonu na Micro:bitu}
    \end{beamercolorbox}%
    \begin{beamercolorbox}[wd=.5\paperwidth,ht=2.5ex,dp=1.125ex,rightskip=6pt,right]{date in head/foot}%
      \usebeamerfont{date in head/foot}\textcolor{gray}{\small\insertframenumber\,/\,\inserttotalframenumber}
    \end{beamercolorbox}}%
}

% --- Česky korektní uvozovky ---
\usepackage{csquotes}
\newcommand{\uv}[1]{\enquote{#1}}

% --- Zdrojový kód (Python) ---
\usepackage{listings}
\lstset{
  language=Python,
  basicstyle=\ttfamily\small,
  backgroundcolor=\color{code-bg},
  frame=single,
  framerule=0pt,
  xleftmargin=6pt,
  xrightmargin=6pt,
  breaklines=true,
  showstringspaces=false,
  keywordstyle=\color{microbit-blue}\bfseries,
  stringstyle=\color{teal},
  commentstyle=\color{gray}\itshape,
  literate=
    {á}{{\'a}}1 {é}{{\'e}}1 {í}{{\'i}}1 {ó}{{\'o}}1 {ú}{{\'u}}1
    {ů}{{\r{u}}}1 {č}{{\v{c}}}1 {š}{{\v{s}}}1 {ž}{{\v{z}}}1
    {Á}{{\'A}}1 {É}{{\'E}}1 {Í}{{\'I}}1 {Ó}{{\'O}}1 {Ú}{{\'U}}1
    {Č}{{\v{C}}}1 {Š}{{\v{S}}}1 {Ž}{{\v{Z}}}1 {ň}{{\v{n}}}1 {ř}{{\v{r}}}1,
}

% --- Zvýraznění pojmů ---
\newcommand{\pojem}[1]{\textbf{\textcolor{microbit-blue}{#1}}}
\newcommand{\pyinline}[1]{\texttt{#1}}

% --- Grafika a diagramy ---
\usepackage{tikz}

% --- Ostatní ---
\usepackage{graphicx}
\usepackage{hyperref}
\hypersetup{colorlinks=true, urlcolor=microbit-blue}


\title{Lekce 03: Podmínky}
\subtitle{Úvod do Pythonu na Micro:bitu}
\date{}

\begin{document}

\begin{frame}
  \titlepage
\end{frame}

% ------------------------------------------------------------------
\begin{frame}{Podmínka v životě}

  Rozhodování podle podmínek děláme pořád:

  \bigskip
  \begin{center}
  \begin{tikzpicture}[node distance=1.2cm]
    \node[fc-start]                        (start) {Jdu ven};
    \node[fc-decision, below=1.4cm of start] (cond)  {Prší?};
    \node[fc-process,  left=2.2cm of cond] (yes)   {Vezmu\\deštník};
    \node[fc-process,  right=2.2cm of cond](no)    {Deštník\\nechám doma};

    \draw[fc-arrow] (start) -- (cond);
    \draw[fc-arrow] (cond)  -- node[above, font=\footnotesize] {Ano} (yes);
    \draw[fc-arrow] (cond)  -- node[above, font=\footnotesize] {Ne}  (no);
  \end{tikzpicture}
  \end{center}
\end{frame}

% ------------------------------------------------------------------
\begin{frame}[fragile]{Podmínka v Pythonu: \texttt{if / else}}

  \begin{lstlisting}
if podminka:
    # co se stane, kdyz je podminka splnena
else:
    # co se stane jinak
  \end{lstlisting}

  \bigskip
  \begin{itemize}
    \item Za \pyinline{if} a \pyinline{else} musí být \textbf{dvojtečka}
    \item Vnitřní příkazy musí být \textbf{odsazeny} (Tab nebo 4 mezery)
    \item Odsazení = součást jazyka — Python ho vyžaduje
  \end{itemize}

  \pause\bigskip
  \begin{alertblock}{Nejčastější chyba}
    Zapomenutá dvojtečka nebo špatné odsazení → \texttt{IndentationError}.\\
    Thonny odsazuje automaticky po dvojtečce — nechej ho pracovat.
  \end{alertblock}
\end{frame}

% ------------------------------------------------------------------
\begin{frame}[fragile]{True a False}

  Podmínka musí být buď \pojem{True} (pravda) nebo \pojem{False} (nepravda).

  \bigskip
  Tlačítko na micro:bitu vrací právě tyto hodnoty:

  \begin{lstlisting}
button_a.is_pressed()   # True, kdyz je stisknuto
                        # False, kdyz neni
  \end{lstlisting}

  \bigskip
  \begin{columns}[T]
    \begin{column}{0.46\textwidth}
      \begin{block}{True (pravda)}
        Podmínka je splněna.\\
        Provede se blok za \pyinline{if}.
      \end{block}
    \end{column}
    \begin{column}{0.46\textwidth}
      \begin{block}{False (nepravda)}
        Podmínka není splněna.\\
        Provede se blok za \pyinline{else}.
      \end{block}
    \end{column}
  \end{columns}
\end{frame}

% ------------------------------------------------------------------
\begin{frame}[fragile]{\texttt{while True:} — program, který běží pořád}

  Bez smyčky by se tlačítko zkontrolovalo jen jednou.\\
  Přidáme \pyinline{while True:} — program se \textbf{opakuje donekonečna}:

  \begin{lstlisting}
from microbit import *

while True:
    if button_a.is_pressed():
        display.show(Image.HEART)
    else:
        display.show(Image.SAD)
  \end{lstlisting}

  \bigskip
  \textcolor{gray}{\small
    \pyinline{while True:} je speciální případ cyklu — podrobně v příští lekci.\\
    Zastavení: tlačítko \textbf{Stop} v Thonny nebo \textbf{Reset} na micro:bitu.
  }
\end{frame}

% ------------------------------------------------------------------
\begin{frame}[fragile]{Více větví: \texttt{if / elif / else}}

  \pyinline{elif} = „jinak jestliže" — přidáme tolik větví, kolik potřebujeme:

  \begin{lstlisting}
while True:
    if button_a.is_pressed():
        display.show(Image.HAPPY)
    elif button_b.is_pressed():
        display.show(Image.ANGRY)
    else:
        display.show(Image.ASLEEP)
  \end{lstlisting}

  \bigskip
  Python testuje podmínky \textbf{odshora dolů} a provede první, která platí.
  Zbytek přeskočí.
\end{frame}

% ------------------------------------------------------------------
\begin{frame}[fragile]{Obě tlačítka najednou: \texttt{and}}

  Klíčové slovo \pojem{and} spojí dvě podmínky — obě musí platit:

  \begin{lstlisting}
while True:
    if button_a.is_pressed() and button_b.is_pressed():
        display.scroll("Obe!")
    elif button_a.is_pressed():
        display.show(Image.ARROW_W)
    elif button_b.is_pressed():
        display.show(Image.ARROW_E)
    else:
        display.show(Image.DIAMOND_SMALL)
  \end{lstlisting}

  \pause\bigskip
  \begin{alertblock}{Proč musí být \texttt{and} jako první?}
    Kdybychom dali \pyinline{and} větev na konec, podmínka \pyinline{button_a.is_pressed()}
    by ji předběhla a nikdy bychom se tam nedostali.
  \end{alertblock}
\end{frame}

% ------------------------------------------------------------------
\begin{frame}{Co jsme se naučili}
  \begin{itemize}
    \item[\checkmark] \pyinline{if podminka:} — provede blok, je-li podmínka \pyinline{True}
    \item[\checkmark] \pyinline{else:} — provede se, když podmínka neplatí
    \item[\checkmark] \pyinline{elif podminka:} — další větev podmínky
    \item[\checkmark] \pyinline{True} / \pyinline{False} — výsledek podmínky
    \item[\checkmark] \pyinline{button_a.is_pressed()} — vrací \pyinline{True} nebo \pyinline{False}
    \item[\checkmark] \pyinline{and} — obě podmínky musí platit zároveň
    \item[\checkmark] \pyinline{while True:} — program běží a kontroluje donekonečna
    \item[\checkmark] Odsazení a dvojtečka jsou povinné — Python je vyžaduje
  \end{itemize}
\end{frame}

\end{document}
